\section*{Sets and Indices}

\[
M = \{1,2,3\}
\]
set of months in the planning horizon, indexed by \(m\).

\section*{Parameters}

For each month \(m \in M\):
\[
p^{\text{buy}}_m
\]
purchase price in month \(m\) (code: \texttt{purchase\_prices[m]}),

\[
p^{\text{sell}}_m
\]
selling price in month \(m\) (code: \texttt{selling\_prices[m]}).
  
Global parameters:
\[
C
\]
warehouse capacity (code: \texttt{warehouse\_capacity}),

\[
I_0
\]
initial stock on January 1 (code: \texttt{initial\_stock}),

\[
F_0
\]
initial funds on January 1 (code: \texttt{initial\_funds}),

\[
I^{\text{end}}
\]
required end-of-quarter inventory level (code: \texttt{end\_inventory}).

\section*{Decision Variables}

For each month \(m \in M\):
\[
b_m \ge 0
\]
quantity purchased in month \(m\) (code: \texttt{buy[m]}),

\[
s_m \ge 0
\]
quantity sold in month \(m\) (code: \texttt{sell[m]}),

\[
I_m \ge 0
\]
stock at the end of month \(m\) (code: \texttt{stock[m]}),

\[
F_m \ge 0
\]
funds at the end of month \(m\) (code: \texttt{funds[m]}).

\section*{Objective}

Maximize final funds net of initial funds:
\[
\max \; F_3 - F_0.
\]

\section*{Constraints}

For month \(m=1\):
\[
I_1 = I_0 - s_1 + b_1,
\]
\[
F_1 = F_0 + p^{\text{sell}}_1 s_1 - p^{\text{buy}}_1 b_1,
\]
\[
I_1 \le C,
\]
\[
s_1 \le I_0.
\]

For months \(m=2,3\):
\[
I_m = I_{m-1} - s_m + b_m,
\]
\[
F_m = F_{m-1} + p^{\text{sell}}_m s_m - p^{\text{buy}}_m b_m,
\]
\[
I_m \le C,
\]
\[
s_m \le I_{m-1}.
\]

End-of-quarter inventory requirement:
\[
I_3 \ge I^{\text{end}}.
\]

Nonnegativity:
\[
b_m \ge 0,\quad s_m \ge 0,\quad I_m \ge 0,\quad F_m \ge 0
\qquad \forall m \in M.
\]

\section*{Notes on Implementation}

The implemented model is a linear program solved by the code with a mathematical programming solver interface.

The stock transition
\[
I_m = I_{m-1} - s_m + b_m
\]
treats purchases in month \(m\) as entering inventory in the same month. The restriction
\[
s_m \le I_{m-1}
\]
enforces the business rule used in the code that grain purchased in month \(m\) cannot be sold in the same month; only stock available before that month's purchase may be sold.

The funds transition
\[
F_m = F_{m-1} + p^{\text{sell}}_m s_m - p^{\text{buy}}_m b_m
\]
combined with \(F_m \ge 0\) implies that end-of-month funds cannot be negative, but there is no separate explicit constraint of the form \(p^{\text{buy}}_m b_m \le F_{m-1}\). Thus purchases may be financed within the month using same-month sales revenue, exactly as represented in the source code.

The objective in the code is \(\max F_3 - F_0\), which is equivalent to maximizing final funds \(F_3\) since \(F_0\) is constant.
